\begin{textsample}{POS Dim 1 – vem\_ed\_35 – Score 52.00 – vem\_ed\_35/t189.txt}  \label{ex:f1_pos_013}
BOAS PRÁTICAS Entre fusos e culturas, a hospitalidade como \textbf{ponto} de \textbf{encontro} Conectar estudantes de diferentes países, culturas e \textbf{rotinas} acadêmicas \textbf{é} um desafio logístico.
Mas, quando o propósito \textbf{é} claro e o \textbf{diálogo} se mantém vivo, a distância se transforma em oportunidade. \textbf{É} nesse território \textbf{que} se desenvolve o Projeto Colaborativo Internacional (PCI) entre as Fatecs e a Setsunan University, do Japão: uma experiência pautada na abordagem COIL (Collaborative Online International Learning) \textbf{que} \textbf{amplia} \textbf{horizontes} e aproxima olhares.
A colaboração \textbf{começou} em 2022 e ganhou escala ao \textbf{longo} dos anos.
No primeiro semestre de 2026, envolveu 30 grupos de alunos de inglês das Fatecs Mogi das Cruzes (professora Neusa Haruka Sezaki Gritti), Zona Leste (professora Solange Bazzon), Garça (professora Leysiane Cristina Pavani Pazini) e Ourinhos (professora Valéria Baccili), evidenciando o potencial de replicabilidade da \textbf{proposta}.
A distribuição dos grupos em várias Fatecs \textbf{foi} uma solução encontrada para equilibrar as equipes mistas de estudantes brasileiros e japoneses, pois o professor Curtis Chu tem cerca de 160 alunos de inglês na Setsunan University, enquanto as \textbf{turmas} de inglês nas Fatecs \textbf{são} \textbf{menores}.
No centro da experiência está uma \textbf{proposta} simples, porém potente.
A \textbf{cada} semana, os estudantes recebem \textbf{tarefas} estruturadas: \textbf{produção} de vídeos, publicações no Padlet e \textbf{interação} com colegas estrangeiros.
A dinâmica \textbf{é} majoritariamente assíncrona, solução necessária diante do fuso \textbf{horário} de 12 horas entre Brasil e Japão.
A \textbf{proposta} vai além do \textbf{intercâmbio} \textbf{linguístico}.
O tema da edição de maio-junho de 2026, hospitalidade, convida os estudantes a refletirem sobre práticas culturais em seus \textbf{contextos}. “Os alunos das Fatecs \textbf{produzem} vídeos sobre a hospitalidade dos brasileiros, e os japoneses fazem a mesma \textbf{coisa} sobre o seu país.
Depois, japoneses comentam como percebem o estilo de hospitalidade brasileira e brasileiros indicam o \textbf{que} sentiram sobre a hospitalidade japonesa”, sintetiza a professora Neusa Haruka Sezaki Gritti, da Fatec Mogi das Cruzes e responsável pela \textbf{coordenação} dos PCIs em inglês na equipe da Apoio à Internacionalização do Ensino Superior da Divisão de Extensão e Pesquisa no Ensino Superior da Coordenadoria Geral de Ensino Superior de Graduação do Centro Paula Souza.
As \textbf{interações} entre os estudantes tornam-se espaço de \textbf{construção} coletiva de sentido. “Eles comentam muito”, \textbf{observa} Gritti, \textbf{destacando} o \textbf{impacto} da atividade na escrita dos alunos. “O professor Curtis leciona segunda-feira de manhã, \textbf{que} \textbf{é} o nosso domingo à noite.
Então, domingo à noite \textbf{começa} a subir um monte de mensagem dos alunos”, \textbf{conta} Pazini.
Por trás do fluxo de atividades, há um intenso trabalho docente.
O planejamento central fica, em grande \textbf{parte}, a cargo do professor Chu. “Ele monta muito bem o plano de PCI e as \textbf{propostas} de \textbf{tarefas}”, \textbf{reconhece} Gritti.
Nas Fatecs, o desafio \textbf{é} manter o engajamento: acompanhar grupos, cobrar entregas e adaptar o projeto à \textbf{realidade} \textbf{local}. \textbf{Diferenças} de calendário acadêmico, prazos curtos e \textbf{níveis} \textbf{variados} de participação impactam o andamento.
Ainda assim, a aposta \textbf{é} no aperfeiçoamento do projeto a \textbf{cada} semestre.
As \textbf{questões} tecnológicas também exigem criatividade.
Muitas atividades \textbf{são} realizadas com recursos mínimos, frequentemente via celular. “Eles gravam no próprio celular e já sobem o vídeo ali”, relata Pazini.
Plataformas como Padlet, Canva, tradutores automáticos e aplicativos de mensagens tornam-se aliados nesse \textbf{processo}.
O engajamento, no \textbf{entanto}, compensa o esforço. “Eu gosto muito de fazer esse PCI”, afirma Pazini, mesmo \textbf{reconhecendo} a carga de trabalho envolvida.
A fala revela um aspecto central do projeto: sua sustentação depende, em grande medida, do compromisso dos docentes.
Esse compromisso também se traduz na parceria.
A relação entre os docentes \textbf{é} marcada por \textbf{diálogo} constante e confiança. “A gente faz de tudo para o negócio fluir”, resume Pazini.
Gritti complementa, entre risos: “Eu sempre cutuco”, referindo-se ao estímulo \textbf{contínuo} necessário para manter o projeto em movimento.
Mais do \textbf{que} um projeto internacional, o PCI se \textbf{consolida} como um laboratório de \textbf{competências} globais.
Ao lidar com \textbf{diferenças} culturais, \textbf{linguísticas} e organizacionais, os estudantes experimentam, na prática, os desafios de um mundo interconectado.
E descobrem \textbf{que}, mesmo entre fusos distantes, há valores \textbf{compartilhados}, como a hospitalidade, capazes de \textbf{criar} \textbf{pontes} duradouras.
Entre ajustes, \textbf{aprendizados} e descobertas, o projeto reafirma uma convicção: quando há colaboração genuína, a distância deixa de \textbf{ser} obstáculo e se transforma em matéria-prima para o \textbf{encontro}.

% matched lemmas: ampliar, aprendizado, cada, coisa, começar, compartilhar, competência, consolidar, construção, contar, contexto, contínuo, coordenação, criar, destacar, diferença, diálogo, encontro, entanto, horizonte, horário, impacto, interação, intercâmbio, linguístico, local, longo, nível, observar, parte, pequeno, ponte, ponto, processo, produzir, produção, proposta, que, questão, realidade, reconhecer, rotina, ser, tarefa, turma, variar
\end{textsample}
